\section{男性盆底肌训练（凯格尔运动·男性版）}

\noindent\textit{【编者注】}盆底肌（pelvic floor muscles）是封闭骨盆下口、承托膀胱、直肠与前列腺的一组肌肉，其中的\textbf{球海绵体肌与坐骨海绵体肌}直接参与阴茎勃起、射精与排尿控制。与女性一样，男性也可以通过"凯格尔运动"训练盆底肌，用于改善\textbf{早泄、勃起硬度、尿失禁与前列腺术后的控尿与性功能}。

\vspace{0.5em}
\noindent\textbf{盆底肌训练对男性的作用}：

\begin{itemize}[leftmargin=2cm]
	\item \textbf{改善早泄}：球海绵体肌参与射精反射，强化盆底肌有助于提高对射精的控制力。多项随机研究显示，盆底肌训练可延长阴道内射精潜伏期（IELT），是早泄行为治疗的重要组成部分
	\item \textbf{改善勃起硬度}：坐骨海绵体肌与球海绵体肌在勃起时压迫阴茎根部静脉、维持海绵体内压力，强化后可改善勃起硬度与维持能力，对轻中度 ED 有辅助作用
	\item \textbf{控尿与术后康复}：前列腺根治术、前列腺增生术（TURP）后常见压力性尿失禁，术前与术后进行盆底肌训练可缩短尿失禁恢复时间
	\item \textbf{慢性前列腺炎/慢性盆腔疼痛（CPPS）}：部分 CPPS 患者盆底肌处于高张力状态，需要的是\textbf{放松}而非单纯强化，应经评估后个体化安排
\end{itemize}

\vspace{0.5em}
\noindent\textbf{如何找到盆底肌}：

\begin{itemize}[leftmargin=2cm]
	\item \textbf{排尿中断法（仅用于辨认，勿常规练习）}：排尿中途尝试夹断尿流，感受用力的肌肉位置
	\item \textbf{憋气法}：想象忍住放屁或"把睾丸上提"的感觉，收缩的即为盆底肌
	\item \textbf{避免的错误}：不要收缩腹部、臀部或大腿；不要屏气；不要把"夹紧"与"肚子用力"混为一谈
\end{itemize}

\vspace{0.5em}
\noindent\textbf{训练方案}：

\begin{enumerate}
	\item \textbf{慢收缩}：缓慢收紧盆底肌，坚持 3--5 秒，再缓慢放松 3--5 秒，重复 10 次
	\item \textbf{快收缩}：快速收紧并立即放松，重复 10 次
	\item \textbf{频率}：每日 3 组，坚持 \textbf{3--6 个月}方可见效；耐力和力量会像其他肌肉一样逐步提升
	\item \textbf{关键原则}：像健身一样循序渐进，感到盆底酸痛应适当休息；不要过度训练（可导致盆底肌疲劳、排尿困难）
\end{enumerate}

\vspace{0.5em}
\noindent\textbf{辅助工具与专业帮助}：

\begin{itemize}[leftmargin=2cm]
	\item 生物反馈、盆底康复仪、电刺激等可帮助找到正确的肌肉并提供反馈，尤其适合"总是练不对"的人
	\item 早泄、ED、尿失禁的患者，建议到泌尿外科或男科、盆底康复门诊接受评估，制定个体化训练方案
\end{itemize}

\begin{tcolorbox}[colback=pink!5!white,colframe=red!50!black,title=贴心小叮咛]
盆底肌训练是"低风险、低成本、可居家进行"的方法，但它需要时间和耐心——多数人需要坚持 3--6 个月才能感受到对射精控制或勃起硬度的改善。如果训练中出现疼痛、排尿困难或症状加重，请及时就医排查盆底肌高张力等问题，而不是继续"加练"。
\end{tcolorbox}
